\section{Parity}
The tests we have made so far in ``if'' statements were about the
orientation of the \robot.  In \python\ it is also
possible to make tests on the values of parameters and of
``for'' loop counters.  It can be useful to test the counter
in the body of a loop in order to perform different actions at
each iteration.  Take for example the drawing of the frieze shaped like a
double comb below:

 {\raggedright
  \includevspacePostScript[1.5cm]{FriseDoublePeigne}{1.0}{1.0}
   }

The teeth of the comb being drawn from left to right, the
$1^{\mathrm{st}}$, $3^{\mathrm{rd}}$,
$5^{\mathrm{th}}$, \ldots\ point downwards whereas
the $2^{\mathrm{nd}}$, $4^{\mathrm{th}}$,
$6^{\mathrm{th}}$, \ldots\ point upwards.  In
the drawing program we use a loop counter {\idt i} which
is 1 when we draw the $1^{\mathrm{st}}$ tooth, 2 when we
draw the $2^{\mathrm{nd}}$, \ldots\ Consequently when the
value of the loop counter is {\em odd} (that is 1, 3, 5,
\ldots) a tooth must be drawn downwards whereas if the value
of the counter is {\em even} (that is 2, 4, 6, \ldots) a tooth must be
drawn upwards.  \index{Parity}A number is odd
when the remainder of its division by 2 is equal to 1. To test that the
value of the counter {\idt i} is odd we therefore use the boolean
expression {\idt i \% 2 == 1} which is {\mcl True} when the remainder
{\idt i \% 2} of the integer division of {\idt i} by 2 is equal to 1 and
{\mcl False} when this remainder is equal to 0, that is when
{\idt i} is even.  Note that the equality test is written with two
equal signs {\idt ==}, the single sign {\idt =} being reserved for the
declaration of constants such as {\idt numberOfTeeth = 30}.

{
\parindent \parindentprog
\parskip 0pt
\nopagebreak\medskip

\carrenoir
\vspace{\neglskip}\mbox{\blc {\mcl from}{\idt \ LibRobot\ }{\mcl import}{\idt \ *}}\par
\nopagebreak
\vspace{0.3\baselineskip}
\vspace{\neglskip}\mbox{\blc {\cmt \#\ Draw\ a\ frieze\ shaped\ like\ a\ double\ comb}}\par
\vspace{\neglskip}\mbox{\blc {\idt numberOfTeeth\ =\ 30}}\par
\vspace{0.3\baselineskip}
\vspace{\neglskip}\mbox{\blc {\idt gl();\ fe()}}\par
\vspace{\neglskip}\mbox{\blc {\mcl for}{\idt \ i\ }{\mcl in}{\idt \ range(1,\ numberOfTeeth\ +\ 1):}}\par
\vspace{\neglskip}\mbox{\blc \blk\blk{\cmt \#\ Horizontal\ line}}\par
\vspace{\neglskip}\mbox{\blc \blk\blk{\idt fd()}}\par
\vspace{\neglskip}\mbox{\blc \blk\blk{\mcl if}{\idt \ i\ \%\ 2\ ==\ 1:}}\par
\vspace{\neglskip}\mbox{\blc \blk\blk\blk\blk{\cmt \#\ Tooth\ pointing\ down}}\par
\vspace{\neglskip}\mbox{\blc \blk\blk\blk\blk{\idt qtr();\ fd();\ ht();\ fd();\ qtr()}}\par
\vspace{\neglskip}\mbox{\blc \blk\blk{\cmt \#\ Tooth\ pointing\ up}}\par
\vspace{\neglskip}\mbox{\blc \blk\blk{\mcl else}{\idt :}}\par
\vspace{\neglskip}\mbox{\blc \blk\blk\blk\blk{\idt qtl();\ fd();\ ht();\ fd();\ qtl()}}\par
\nopagebreak
\vspace{\neglskip}\mbox{\blc {\idt st()}}\par

\nopagebreak\medskip }\nopagebreak


\section{Comparison of integers}

If {\idt i} is a loop counter or an integer parameter of a procedure,
the boolean expression {\idt i > 10} is {\mcl True}
if the value of {\idt i} is strictly greater than 10 and {\mcl
False} otherwise, that is if the value of {\idt
i} is less than or equal to 10.  The mathematical symbols $\leq$
(less than or equal) and $\geq$ (greater than or equal) are written {\idt <=} and
{\idt >=} in \python\ (because the mathematical symbols $\leq$ and $\geq$ are
not found on computer keyboards whereas {\idt <}, {\idt >}
and {\idt =} are always there).  For example the boolean
expressions {\idt 10 <= 10} and {\idt 10 <= 20} are {\mcl True}
whereas {\idt 10 <= 9} is {\mcl False}.  In
\python, as we have seen, equality is written with two equal signs {\idt ==}
since the usual mathematical symbol {\idt =} is used to declare constants.
The mathematical symbol $\neq$ (different) is written {\ttfamily\idt !=}.
For example {\idt 10 == 10} and {\ttfamily\idt 10 != 20} are {\mcl True}
whereas {\ttfamily\idt 10 != 10} and {\idt 10 == 20} are {\mcl
False}\footnote{\ If by mistake one writes {\mcl if}\ {\idt i = 10:}
instead of {\mcl if}\ {\idt i == 10:}, \python\ rejects the program with
the message {\idt SyntaxError: invalid syntax. Maybe you meant '==' or ':='
instead of '='?}}.
Boolean expressions comparing
integers or values of integer parameters can be used in the test
of the ``if'' statement.  For example the command {\idt fdn(n)} which
makes the \robot\ move forward {\idt n} times if the value of {\idt n} is
positive, makes it move backward if the value of {\idt n} is negative and which
does not change its position when the value of {\idt n} is zero can
be defined as follows:

\input python/DefinitionFdn.tex

\noindent Note that the negative actual parameter of the call
{\idt fdn(-5)} is simply written with its sign, between the parentheses of
the call. To move backward, the \robot\ makes a half turn, moves forward {\idt -n} times
(which is positive) and makes a half turn again.

More generally, we can use the operators {\idt <} (strictly less than), {\idt <=} (less than or equal),
{\idt ==} (equal), {\ttfamily\idt !=} (different), {\idt >=} (greater than or equal) and {\idt >} (strictly greater than) to compare the values
of two integer expressions, as for example {\idt i <= l} in the following program which draws the paper streamer
below. The command {\idt gs} accepts an integer, as
in {\idt gs(20)} or {\idt gs(i)}, as well as a number with a decimal part, written
with a point as in {\idt gs(20.5)} (we will come back to this
in chapter 19, devoted to rational expressions):

 \includevspacePostScript[1.5cm]{Serpentin}{0.95}{0.95}
 \input python/Streamer.tex

\section{Negation, conjunction and disjunction}

The {\em negation} \index{Negation} of a boolean expression $E$ is written {\mcl not} $E$, which reads `not $E$'. If the
value of $E$ is {\mcl True} then that of {\mcl not} $E$ is {\mcl False} whereas if the value of $E$ is {\mcl False} then
that of {\mcl not} $E$ is {\mcl True}. For example  {\mcl not} {\idt (5 \% 2 == 1)} is {\mcl False} since
{\idt 5 \% 2 == 1} is {\mcl True}, {\idt 5} being odd. The parentheses are
not indispensable but make the expression more readable (we will see why in the Addenda). This can be summarised by the following truth table:

\begin{center}
\begin{tabular}{|c|c|} \hline
$E$ & {\mcl not} $E$\\ \hline \hline
{\mcl True} & {\mcl False}\\ \hline
{\mcl False} & {\mcl True}\\ \hline
\end{tabular}
\end{center}

The {\em conjunction} \index{Conjunction} of two boolean expressions $E_{1}$ and $E_{2}$ is written $E_{1}$ {\mcl and} $E_{2}$,
which reads `$E_{1}$ and $E_{2}$'. The value of $E_{1}$ {\mcl and} $E_{2}$ is {\mcl True} if and
only if the value of $E_{1}$ and that of $E_{2}$ are {\mcl True}. If either of the boolean
expressions $E_{1}$ or
$E_{2}$ is {\mcl False} then the conjunction $E_{1}$ {\mcl and} $E_{2}$ is {\mcl False} too.
For example  {\idt 5 >= 0} {\mcl and} {\idt 5 \% 2 == 1} is {\mcl True} since
{\idt 5} is positive and odd. This can be summarised by the following truth table:

\begin{center}
\begin{tabular}{|c|c|c|} \hline
$E_{1}$ & $E_{2}$ & $E_{1}$ {\mcl and} $E_{2}$\\ \hline \hline
{\mcl True} & {\mcl True} & {\mcl True}\\ \hline
{\mcl True} & {\mcl False} & {\mcl False}\\ \hline
{\mcl False} & {\mcl True} & {\mcl False}\\ \hline
{\mcl False} & {\mcl False} & {\mcl False}\\ \hline
\end{tabular}
\end{center}

The {\em disjunction} \index{Disjunction} of two boolean expressions $E_{1}$ and $E_{2}$ is written $E_{1}$ {\mcl or} $E_{2}$,
which reads `$E_{1}$ or $E_{2}$'. The value of $E_{1}$ {\mcl or} $E_{2}$ is {\mcl True} if and
only if the value of $E_{1}$ or that of $E_{2}$ is {\mcl True}. If the values of $E_{1}$ and of
$E_{2}$ are  {\mcl False} then the disjunction $E_{1}$ {\mcl or} $E_{2}$ is {\mcl False} too.
For example  {\idt 4 >= 0} {\mcl or} {\idt 4 \% 2 == 1} is {\mcl True} since
{\idt 4} is positive, even though {\idt 4} is even. This can be summarised by the following truth table:

\begin{center}
\begin{tabular}{|c|c|c|} \hline
$E_{1}$ & $E_{2}$ & $E_{1}$ {\mcl or} $E_{2}$\\ \hline \hline
{\mcl True} & {\mcl True} & {\mcl True}\\ \hline
{\mcl True} & {\mcl False} & {\mcl True}\\ \hline
{\mcl False} & {\mcl True} & {\mcl True}\\ \hline
{\mcl False} & {\mcl False} & {\mcl False}\\ \hline
\end{tabular}
\end{center}

Conjunction and disjunction can be implemented simply with an electric circuit as shown in figure
\ref{fig:andor}
below
%(page \pageref{fig:andor})
with the con\-ven\-tion that {\mcl True} represents a closed switch
and a lit bulb (the current flows) and {\mcl False} represents an open switch
and an unlit bulb (the current does not flow).

\begin{figure}[htbp]
 \includevspacePostScript[4.3cm]{CircuitsAndOr}{0.95}{0.95}
 \caption{Electric circuits implementing logical conjunction and\label{fig:andor}}
\centering{disjunction}
\end{figure}

A disjunction is used in the following program to test whether the \robot\ is facing north or
east:

\vspace{0.3cm}
 {\raggedright
 \includevspacePostScript[0.5cm]{FriseEnOr}{1.0}{1.0}
{
\parindent \parindentprog
\parskip 0pt
\nopagebreak\medskip

\carrenoir
\vspace{\neglskip}\mbox{\blc {\mcl from}{\idt \ LibRobot\ }{\mcl import}{\idt \ *}}\par
\nopagebreak
\vspace{0.3\baselineskip}
\vspace{\neglskip}\mbox{\blc {\cmt \#\ Draw\ a\ golden\ frieze}}\par
\nopagebreak
\vspace{\neglskip}\mbox{\blc {\idt length\ =\ 31}}\par
\nopagebreak
\vspace{0.3\baselineskip}
\vspace{\neglskip}\mbox{\blc {\idt gl()}}\par
\nopagebreak
\vspace{\neglskip}\mbox{\blc {\mcl for}{\idt \ i\ }{\mcl in}{\idt \ range(length):}}\par
\nopagebreak
\vspace{\neglskip}\mbox{\blc \blk\blk{\mcl if}{\idt \ tn()\ }{\mcl or}{\idt \ te():}}\par
\nopagebreak
\vspace{\neglskip}\mbox{\blc \blk\blk\blk\blk{\idt cr()}}\par
\nopagebreak
\vspace{\neglskip}\mbox{\blc \blk\blk{\mcl else}{\idt :}}\par
\nopagebreak
\vspace{\neglskip}\mbox{\blc \blk\blk\blk\blk{\idt fdn({-}1);\ cl();\ qtl()}}\par
\nopagebreak
\vspace{\neglskip}\mbox{\blc {\idt st()}}\par

\nopagebreak\medskip }\nopagebreak

  }

\section{Boolean expressions}

The table below summarises the rules for writing boolean expressions and the rules for computing their value
in \python:

\index{Boolean expression} \index{Test}

\begin{center}
\protect\frenchspacing
\begin{tabular}{@{}|l|p{6.8cm}|@{}} \hline
\multicolumn{2}{|l|}{\em Boolean expressions:} \\ \hline \hline
{\mcl True} & Value `true'.\\
{\mcl False} & Value `false'.\\ \hline
{\small\rm Integer expr.$_{1}$} {\idt <=} {\small\rm Integer expr.$_{2}$}& {\mcl True} if and only if the value of {\small\rm Integer expr.$_{1}$} is
                                 less than or equal to that of {\small\rm Integer expr.$_{2}$}.\\
{\small\rm Integer expr.$_{1}$} {\idt <} {\small\rm Integer expr.$_{2}$}& {\mcl True} if and on\-ly if the va\-lue of {\small\rm Integer expr.$_{1}$} is
                                 strictly less than that of {\small\rm Integer expr.$_{2}$}.\\
{\small\rm Integer expr.$_{1}$} {\idt ==} {\small\rm Integer expr.$_{2}$}& {\mcl True} if and on\-ly if the value of {\small\rm Integer expr.$_{1}$} is
                                 equal to that of {\small\rm Integer expr.$_{2}$}.\\
{\small\rm Integer expr.$_{1}$} {\ttfamily\idt !=} {\small\rm Integer expr.$_{2}$}& {\mcl True} if and on\-ly if the value of {\small\rm Integer expr.$_{1}$} is
                                 different from that of {\small\rm Integer expr.$_{2}$}.\\ \hline
\end{tabular}
\end{center}
\begin{center}
\protect\frenchspacing
\begin{tabular}{@{}|l|p{6.8cm}|@{}} \hline
{\small\rm Integer expr.$_{1}$} {\idt >} {\small\rm Integer expr.$_{2}$}& {\mcl True} if and on\-ly if the va\-lue of {\small\rm Integer expr.$_{1}$} is
                                 strictly greater than that of {\small\rm Integer expr.$_{2}$}.\\
{\small\rm Integer expr.$_{1}$} {\idt >=} {\small\rm Integer expr.$_{2}$}& {\mcl True} if and only if the value of {\small\rm Integer expr.$_{1}$} is
                                greater than or equal to that of {\small\rm Integer expr.$_{2}$}.\\ \hline
                                \end{tabular}
\end{center}
\begin{center}
\protect\frenchspacing
\begin{tabular}{@{}|l|p{6.8cm}|@{}} \hline
{\mcl not} {\small\rm Bool. expr.} & {\mcl True} if and only if the value of {\small\rm Bool. expr.} is
                                equal to {\mcl False}.\\ \hline
{\small\rm Bool. expr.$_{1}$} {\mcl and} {\small\rm Bool. expr.$_{2}$}& {\mcl True} if and only if the values of {\small\rm Bool. expr.$_{1}$} and
                                 of {\small\rm Bool. expr.$_{2}$} are equal to {\mcl True}.\\
{\small\rm Bool. expr.$_{1}$} {\mcl or} {\small\rm Bool. expr.$_{2}$}& {\mcl True} if and only if the value of {\small\rm Bool. expr.$_{1}$}
                                 or that of {\small\rm Bool. expr.$_{2}$} is equal to {\mcl True}.\\ \hline
\end{tabular}
\end{center}

\index{Boolean function}A boolean expression can be given a name to
make programs more readable. For example:

\noindent\begin{minipage}{\textwidth}
{
\parindent \parindentprog
\medskip

\vspace{\neglskip}\mbox{\blc {\cmt \#\ is\ n\ odd?}}\par
\vspace{\neglskip}\mbox{\blc {\mcl def}{\idt \ isOdd(n):}}\par
\vspace{\neglskip}\mbox{\blc \blk\blk{\mcl return}{\idt \ n\ \%\ 2\ ==\ 1}}\par
\vspace{0.3\baselineskip}
\vspace{\neglskip}\mbox{\blc {\cmt \#\ is\ n\ even?}}\par
\vspace{\neglskip}\mbox{\blc {\mcl def}{\idt \ isEven(n):}}\par
\vspace{\neglskip}\mbox{\blc \blk\blk{\mcl return}{\idt \ n\ \%\ 2\ ==\ 0}}\par

\medskip }
\end{minipage}

\noindent These declarations look like declarations of procedures with parameters
but their body is not a sequence of instructions for the
\robot: the {\mcl return} instruction is followed by a
boolean expression whose value is computed at each call and returned as
the result of the call.
We say that {\idt isOdd} and {\idt isEven} are {\em boolean functions} because their evaluation always returns a boolean (a value of type {\idt bool}, that is {\mcl True} or {\mcl False}).
For example {\idt isOdd(5)} is {\mcl True} and {\idt isEven(5)} is
{\mcl False}. They can be used in a test such as {\mcl if}\ {\idt isOdd(i):}
\ldots\ or in a boolean expression such as {\mcl not}\ {\idt isOdd(i)}. The orientation
tests {\idt tn()}, {\idt te()}, \ldots\ of the \robot\ are also boolean
functions, declared in the {\idt LibRobot} library.

\needspace{6\baselineskip}\noindent \rule{137mm}{0.1mm}\par\nopagebreak %---------------------------------------------
\begin{exercices}{Write \python\ programs using parity tests to produce
the following drawings:

\tiret  the pyramidal serpentine and the grid of figure \ref{fig:serpyrgricr} below;
%(page \pageref{fig:serpyrgricr});

\tiret   the octagonal tiling and the tiling of interlaced hexagons of figure \ref{fig:pavocthex} (page \pageref{fig:pavocthex}).

\noindent $\Box$}\end{exercices}
\begin{figure}[htbp]
  \includevspacePostScript[6.0cm]{SerpPyr+GrilleDeCr}{0.95}{0.95}
 \caption{Pyramidal serpentine and grid of crosses}\label{fig:serpyrgricr}
\end{figure}
\begin{figure}[htbp]
\includevspacePostScript[5.4cm]{PavageOctoHexa}{0.95}{0.95}
 \caption{Tilings of juxtaposed octagons and of interlaced hexagons}\label{fig:pavocthex}
\end{figure}
\begin{exercices}{Write \python\ programs using boo\-le\-an expressions comparing
integer values to produce the following drawings (the second one is difficult):

\tiret  the digital digits of figure \ref{fig:alphabet} (page \pageref{fig:alphabet}) where the digit to draw
is passed as a parameter of the drawing procedure;

\tiret  the tiling of cogwheels of figure \ref{fig:pavrd} below;

\tiret  the Egyptian carpet of figure \ref{fig:tapegyp} (page \pageref{fig:tapegyp}).

\noindent $\Box$}\end{exercices}
\begin{figure}[htbp]
 \includevspacePostScript[5.7cm]{PavageRoueDentee}{0.95}{0.95}
 \caption{Tiling of cogwheels}\label{fig:pavrd}
\end{figure}
\begin{figure}[htbp]
 \includevspacePostScript[7.2cm]{TapisEgyptien}{0.95}{0.95}
 \caption{Egyptian carpet\label{fig:tapegyp}}
\end{figure}
\begin{exercices}{Write \python\ programs using negations, conjunctions or disjunctions of boolean
expressions to produce the following drawings (the second one is difficult):

\tiret  the zebra-striped square of figure \ref{fig:carzeb} (page \pageref{fig:carzeb});

\tiret  the Alhambra mosaic shown in figure \ref{fig:mosaic6} (page \pageref{fig:mosaic6}).

\noindent $\Box$}\end{exercices}
\begin{figure}[htbp]
 \includevspacePostScript[5.8cm]{CarreZebree}{0.95}{0.95}
 \caption{Zebra-striped square}\label{fig:carzeb}
\end{figure}
\begin{figure}[htbp]
 \includevspacePostScript[5.7cm]{MosaiqueAlhambra6}{0.95}{0.95}
 \caption{Alhambra mosaic}\label{fig:mosaic6}
\end{figure}
%%%%%%%%%%%%%%%%%%%%%%%%%%%%%%%%%%%%%%%%%%%%%%%%%%%%%%%%%%
\begin{corrige}{
We give one program, the others being in the folder
{\idt book/\allowbreak 16\_boolean\_expressions}:
\input python/PyramidalSerpentine.tex
\noindent $\Box$}\end{corrige}
\begin{corrige}{
We give one program, the third one being in the folder
{\idt book/\allowbreak 16\_boolean\_expressions}:
\input python/CogwheelsTiling.tex
\noindent $\Box$}\end{corrige}
\begin{corrige}{
We give one program, the first one being in the folder
{\idt book/\allowbreak 16\_boolean\_expressions}:
\input python/AlhambraMosaic6.tex
\noindent $\Box$}\end{corrige}

%%%%%%%%%%%%%%%%%%%%%%%%%%%%%%%%%%%%%%%%%%%%%%%%%%%%%%%%%%%%%%%%%%%%%%%%%%%
\noindent \framebox[\textwidth]{
\begin{minipage}{\myparboxwidth}
\setlength{\parindent}{2em}
\small \sf
\vspace{\abovedisplayskip}
\section*{Addenda}

 In \python\ the priorities of integer and boolean operators are defined as follows:

 \vspace{\baselineskip}

 \index{Operator priority}
\begin{center}
\begin{tabular}{|l|l|l|} \hline
\multicolumn{3}{|l|}{\em Priority of integer and boolean operators:} \\ \hline \hline
{\idt f(x)} (call) & function calls & highest priority \\ \hline
$-$ (unary) & opposite & \\ \hline
$\ast$, {\idt //}, {\idt \%} & multiplicative operators & \\ \hline
$+$, $-$ (binary) & additive operators & \\ \hline
{\idt <}, {\idt <=}, {\idt ==}, {\ttfamily\idt !=}, {\idt >=}, {\idt >}  & relational operators & \\ \hline
{\mcl not} & negation & \\ \hline
{\mcl and} & conjunction & \\ \hline
{\mcl or} & disjunction & lowest priority \\ \hline
\end{tabular}
\end{center}

\vspace{\baselineskip}

 \noindent For example the boolean expression `{\mcl True\ and\ False\ or\ False}' is eva\-lu\-at\-ed as
 `{\idt (}{\mcl True\ and\ False}{\idt )}\ {\mcl or\ False}' and is therefore equal to {\mcl False}.

\vspace{\abovedisplayskip}
\end{minipage}
}
%%%%%%%%%%%%%%%%%%%%%%%%%%%%%%%%%%%%%%%%%%%%%%%%%%%%%%%%%%%%%%%%%%%%%%%%%%%%

%%%%%%%%%%%%%%%%%%%%%%%%%%%%%%%%%%%%%%%%%%%%%%%%%%%%%%%%%%%%%%%%%%%%%%%%%%%
\noindent \framebox[\textwidth]{
\begin{minipage}{\myparboxwidth}
\setlength{\parindent}{2em}
\small \sf
\vspace{\abovedisplayskip}
\indent These priorities are close to the usual mathematical conventions. For
example `{\idt 0\ <=\ i}\ {\mcl and}\ {\idt i\ <\ 100}' is evaluated as
`{\idt (0\ <=\ i)}\ {\mcl and}\ {\idt (i\ <\ 100)}', which is indeed what we mean. This was not
the case in older languages such as \pascal, where this expression was
incorrect and all the parentheses had to be written. \python\ even allows
writing `{\idt 0\ <=\ i\ <\ 100}' as in mathematics.

Since {\mcl not} has a lower priority than the relational operators,
`{\mcl not}\ {\idt i\ \%\ 2\ ==\ 1}' is evaluated as `{\mcl not}\ {\idt (i\ \%\ 2\ ==\ 1)}',
which is indeed what we mean. This is not the case in other languages
such as \pascal\ or \ocaml, where {\mcl not} applies to what immediately follows it,
like a function call, so that the parentheses must be written.
When in doubt, one can always add parentheses.

Finally, \python\ evaluates conjunctions and disjunctions from left to right and stops as soon
as the result is known: in $E_{1}$ {\mcl and} $E_{2}$, if $E_{1}$ is {\mcl False}
then $E_{2}$ is not computed since the conjunction is necessarily {\mcl False}; likewise
in $E_{1}$ {\mcl or} $E_{2}$, if $E_{1}$ is {\mcl True} then $E_{2}$ is not
computed.

The term ``boolean'' ({\idt bool} in \python) comes from Boole\footnote{\ \href{https://en.wikipedia.org/wiki/George_Boole}{George Boole}, British mathematician, 1815--1864.},
 creator of modern mathematical logic.

\vspace{\abovedisplayskip}
\end{minipage}
}
%%%%%%%%%%%%%%%%%%%%%%%%%%%%%%%%%%%%%%%%%%%%%%%%%%%%%%%%%%%%%%%%%%%%%%%%%%%%
\newpage                 % If the chapter ends with a blank page
\thispagestyle{empty}
